% Source PDF 1–5. Existing Chinese source text transcribed, not retranslated.
% English copyright notice translated below; source titles/person names preserved.
\definecolor{fgablue}{RGB}{0,36,83}
\definecolor{fgapaper}{RGB}{245,232,195}
% Native vector reconstruction of the cover's temple/circular motif.
% No original publisher raster logo or scanned page is embedded.
\newcommand{\fgatemple}[1]{%
\begin{tikzpicture}[x=1mm,y=1mm,scale=#1]
\fill[fgablue] (-47,-34) rectangle (47,-30);
\fill[fgablue] (-44,-28) rectangle (44,-24);
\fill[fgablue] (-41,-22) rectangle (41,-18);
\fill[fgablue] (-38,-16) rectangle (38,-12);
\foreach \x in {-32,-11,10,31}{
 \fill[fgablue] (\x-4,-10)--(\x+5,-10)--(\x+4,24)--(\x-3,24)--cycle;
 \fill[fgablue,rounded corners=2mm] (\x-5,26) rectangle (\x+6,30);}
\fill[fgablue] (-39,32) rectangle (40,36);
\fill[fgablue] (-41,38) rectangle (42,41);
\fill[fgablue] (-43,42)--(0,62)--(44,42)--(42,38)--(0,57)--(-41,38)--cycle;
\fill[fgablue] (-33,43)--(0,56)--(34,43)--cycle;
\draw[fgablue,line width=6mm] (-60,-9) arc[start angle=198,end angle=405,radius=64mm];
\end{tikzpicture}}

\begingroup
\let\cleardoublepage\clearpage
\begin{titlepage}
\begin{tikzpicture}[remember picture,overlay]
\fill[fgapaper] (current page.south west) rectangle (current page.north east);
\node[anchor=north west,text width=118mm,align=left] at ([xshift=51mm,yshift=-11mm]current page.north west)
 {\sffamily\bfseries\fontsize{27}{32}\selectfont 基础代数几何\\[3mm]
 \fontsize{13}{17}\selectfont Fundamental Algebraic Geometry\\[2mm]
 \fontsize{12}{16}\selectfont Grothendieck's FGA Explained\\[5mm]
 Grothendieck\\《基础代数几何学（FGA）》解读\\[6mm]
 \normalfont\sffamily\fontsize{11}{15}\selectfont
 Barbara Fantechi, Lothar Göttsche\\
 Luc Illusie, Steven L. Kleiman\\Nitin Nitsure, Angelo Vistoli};
\node[anchor=north west,align=left] at ([xshift=14mm,yshift=-11mm]current page.north west)
 {\sffamily\fontsize{7}{9}\selectfont 美国数学会经典影印系列\\\fontsize{23}{27}\selectfont AMS};
\node at ([xshift=2mm,yshift=-46mm]current page.center){\fgatemple{0.92}};
\node[anchor=south west] at ([xshift=16mm,yshift=14mm]current page.south west){\large 高等教育出版社};
\end{tikzpicture}
\null
\end{titlepage}

% Source PDF 2: back-cover copy, placed in the scan's sequence.
\begin{titlepage}
\begin{tikzpicture}[remember picture,overlay]
\fill[fgapaper] (current page.south west) rectangle (current page.north east);
\end{tikzpicture}
\vspace*{-7mm}
\noindent\begin{minipage}[t]{104mm}\small\sffamily
\setlength{\parindent}{1.5em}
Alexander Grothendieck 以极其深刻、极富创造性的思想, 使得代数几何学发生了里程碑式的变革. 他在 1957 年到 1962 年的布尔巴基讨论班上给出了他的新理论的一个概述, 然后将这些讲义整理成一系列的文章, 编成了著名的《基础代数几何学》（Fondements de la géométrie algébrique）, 即我们熟知的 FGA.

FGA 中的许多内容目前已广为人知, 然而仍有一些知识是大家所不了解的, 只有少数几何学家熟悉它的全部内容. 本书源自 2003 年在意大利的里雅斯特（Trieste）开设的基础代数几何高级学校, 目的就是完善 Grothendieck 对于其理论过于简要的概述. 本书讨论的四个重要主题为:

下降理论、Hilbert 和 Quot 概形、形式存在定理和 Picard 概形. 作者们给出了主要结果的完整证明, 在必要时使用较新的概念以使读者更好理解, 并且阐述了 FGA 的理论与新近发展的联系.

本书适合于对代数几何学感兴趣的研究生和专业研究人员阅读. 学习本书需要全面扎实的基础概形理论知识.
\end{minipage}\hfill
\begin{minipage}[t]{29mm}\fontsize{7}{10}\selectfont\sffamily
本版只限于中华人民共和国境内发行. 本版经由美国数学会授权仅在中华人民共和国境内销售, 不得出口.
\end{minipage}
\par\vspace{12mm}\noindent\sffamily 美国数学会经典影印系列\\{\Huge AMS}
\vfill
\begin{flushright}\ttfamily ISBN 978-7-04-051012-6\\
% Native EAN-13 bars for source ISBN 9787040510126.
\begin{tikzpicture}[x=.25mm,y=1mm]
\fill (0,0) rectangle (1,14);
\fill (2,0) rectangle (3,14);
\fill (4,0) rectangle (5,14);
\fill (5,0) rectangle (6,14);
\fill (6,0) rectangle (7,14);
\fill (8,0) rectangle (9,14);
\fill (9,0) rectangle (10,14);
\fill (13,0) rectangle (14,14);
\fill (16,0) rectangle (17,14);
\fill (19,0) rectangle (20,14);
\fill (23,0) rectangle (24,14);
\fill (27,0) rectangle (28,14);
\fill (28,0) rectangle (29,14);
\fill (30,0) rectangle (31,14);
\fill (33,0) rectangle (34,14);
\fill (34,0) rectangle (35,14);
\fill (35,0) rectangle (36,14);
\fill (37,0) rectangle (38,14);
\fill (41,0) rectangle (42,14);
\fill (42,0) rectangle (43,14);
\fill (44,0) rectangle (45,14);
\fill (46,0) rectangle (47,14);
\fill (48,0) rectangle (49,14);
\fill (50,0) rectangle (51,14);
\fill (53,0) rectangle (54,14);
\fill (54,0) rectangle (55,14);
\fill (55,0) rectangle (56,14);
\fill (57,0) rectangle (58,14);
\fill (58,0) rectangle (59,14);
\fill (61,0) rectangle (62,14);
\fill (62,0) rectangle (63,14);
\fill (64,0) rectangle (65,14);
\fill (65,0) rectangle (66,14);
\fill (66,0) rectangle (67,14);
\fill (69,0) rectangle (70,14);
\fill (71,0) rectangle (72,14);
\fill (72,0) rectangle (73,14);
\fill (75,0) rectangle (76,14);
\fill (76,0) rectangle (77,14);
\fill (78,0) rectangle (79,14);
\fill (79,0) rectangle (80,14);
\fill (81,0) rectangle (82,14);
\fill (82,0) rectangle (83,14);
\fill (85,0) rectangle (86,14);
\fill (87,0) rectangle (88,14);
\fill (92,0) rectangle (93,14);
\fill (94,0) rectangle (95,14);
\end{tikzpicture}
\\\normalfont 定价 135.00 元\end{flushright}
\end{titlepage}

% Source PDF 3: original title leaf, Chinese-adapted type hierarchy.
\begin{titlepage}
\vspace*{4mm}
\noindent\begin{minipage}[t]{35mm}\fontsize{7}{10}\selectfont 美国数学会经典影印系列\\{\Huge\sffamily AMS}\end{minipage}%
\begin{minipage}[t]{106mm}\raggedright\sffamily\bfseries
{\fontsize{24}{30}\selectfont 基础代数几何\\}
\vspace{3mm}{\Large Fundamental Algebraic Geometry\\}
\vspace{3mm}{\large Grothendieck's FGA Explained\\}
\vspace{8mm}{\large Grothendieck《基础代数几何学（FGA）》解读\\}
\vspace{8mm}\normalfont\sffamily\large
Barbara Fantechi, Lothar Göttsche\\
Luc Illusie, Steven L. Kleiman\\
Nitin Nitsure, Angelo Vistoli
\end{minipage}
\vspace{12mm}\begin{center}\fgatemple{0.70}\end{center}
\vfill\begin{center}\Large 高等教育出版社\,·\,北京\end{center}
\end{titlepage}

% Source PDF 4: all bibliographic, copyright, CIP and production matter.
\begin{titlepage}\fontsize{8}{11}\selectfont
\noindent 图字: 01-2018-2924 号
\par\vspace{5mm}
\textit{Fundamental Algebraic Geometry: Grothendieck's FGA Explained}, 作者 Barbara Fantechi、Lothar Göttsche、Luc Illusie、Steven L. Kleiman、Nitin Nitsure 和 Angelo Vistoli, 最初由美国数学会出版. 版权 \copyright\ 2005 年美国数学会所有. 保留一切权利. 本影印版由高等教育出版社有限公司经美国数学会授权并许可出版. 本特别版仅供中华人民共和国境内发行. 美国数学会授权本版仅在中华人民共和国境内销售, 不得出口.

本书原版最初由美国数学会于 2005 年出版, 原书名为 \textit{Fundamental Algebraic Geometry: Grothendieck's FGA Explained}, 作者为 Barbara Fantechi, Lothar Göttsche, Luc Illusie, Steven L. Kleiman, Nitin Nitsure, Angelo Vistoli. 美国数学会保留原书所有版权.

原书版权声明: Copyright \copyright\ 2005 by the American Mathematical Society.

本影印版由高等教育出版社有限公司经美国数学会独家授权出版.

本版只限于中华人民共和国境内发行. 本版经由美国数学会授权仅在中华人民共和国境内销售, 不得出口.
\par\vspace{6mm}
\noindent\begin{minipage}[t]{43mm}\bfseries Grothendieck\\《基础代数几何学\\（FGA）》解读\\\normalfont\sffamily
Grothendieck Jichu\\Daishu Jihexue (FGA) Jiedu
\end{minipage}\hfill
\begin{minipage}[t]{91mm}
\textbf{图书在版编目（CIP）数据}

Grothendieck《基础代数几何学（FGA）》解读 = Fundamental Algebraic Geometry: Grothendieck's FGA Explained: 英文 /（意）凡泰基（Barbara Fantechi）等著. — 影印本. — 北京: 高等教育出版社, 2019.1

ISBN 978-7-04-051012-6

I. {\fontspec{Noto Serif CJK SC}①}G…\quad II. {\fontspec{Noto Serif CJK SC}①}凡…\quad III. {\fontspec{Noto Serif CJK SC}①}代数几何—教材—英文\quad IV. {\fontspec{Noto Serif CJK SC}①}O187

中国版本图书馆 CIP 数据核字（2018）第 265631 号
\end{minipage}
\par\vfill
\noindent 策划编辑\quad 李\quad 鹏\qquad 责任编辑\quad 李\quad 鹏\\
封面设计\quad 张申申\qquad 责任印制\quad 陈伟光
\par\vspace{5mm}
\noindent\begin{minipage}[t]{69mm}\fontsize{7}{10}\selectfont
出版发行\quad 高等教育出版社\\
社址\quad 北京市西城区德外大街4号\\
邮政编码\quad 100120\\
购书热线\quad 010-58581118\\
咨询电话\quad 400-810-0598\\
网址\quad http://www.hep.edu.cn\\
http://www.hep.com.cn\\
网上订购\quad http://www.hepmall.com.cn\\
http://www.hepmall.com\\
http://www.hepmall.cn\\
印刷\quad 北京新华印刷有限公司
\end{minipage}\hfill
\begin{minipage}[t]{66mm}\fontsize{7}{10}\selectfont
开本\quad $787\,\mathrm{mm}\times1092\,\mathrm{mm}$\quad 1/16\\
印张\quad 22.5\\
字数\quad 540 千字\\
版次\quad 2019 年 1 月第 1 版\\
印次\quad 2019 年 1 月第 1 次印刷\\
定价\quad 135.00 元\\[3mm]
本书如有缺页、倒页、脱页等质量问题, 请到所购图书销售部门联系调换.\\
版权所有\quad 侵权必究\\
［物料号 51012-00］
\end{minipage}
\end{titlepage}

% Source PDF 5. Existing Chinese publisher's foreword faithfully transcribed.
\chapter*{出版者的话}
近年来, 我国的科学技术取得了长足进步, 特别是在数学等自然科学基础领域不断涌现出一流的研究成果. 与此同时, 国内的科研队伍与国外的交流合作也越来越密切, 越来越多的科研工作者可以熟练地阅读英文文献, 并在国际顶级期刊发表英文学术文章、在国外出版社出版英文学术著作.

然而, 在国内阅读海外原版英文图书仍不是非常便捷. 一方面, 这些原版图书主要集中在科技、教育比较发达的大中城市的大型综合图书馆以及科研院所的资料室中, 普通读者借阅不甚容易; 另一方面, 原版书价格昂贵, 动辄上百美元, 购买也很不方便. 这极大地限制了科技工作者对于国外先进科学技术知识的获取, 间接阻碍了我国科技的发展.

高等教育出版社本着植根教育、弘扬学术的宗旨服务我国广大科技和教育工作者, 同美国数学会（American Mathematical Society）合作, 在征求海内外众多专家学者意见的基础上, 精选该学会近年出版的数十种专业著作, 组织出版了“美国数学会经典影印系列”丛书. 美国数学会创建于 1888 年, 是国际上极具影响力的专业学术组织, 目前拥有近 30000 会员和 580 余个机构成员, 出版图书 3500 多种, 冯·诺依曼、莱夫谢茨、陶哲轩等世界级数学大家都是其作者. 本影印系列涵盖了代数、几何、分析、方程、拓扑、概率、动力系统等所有主要数学分支以及新近发展的数学主题.

我们希望这套书的出版, 能够对国内的科研工作者、教育工作者以及青年学生起到重要的学术引领作用, 也希望今后能有更多的海外优秀英文著作被介绍到中国.
\begin{flushright}高等教育出版社\\2016 年 12 月\end{flushright}
\clearpage
\endgroup
